\documentclass[12pt,a4paper]{article}

\usepackage[utf8]{inputenc}
\usepackage[T1]{fontenc}
\usepackage[spanish]{babel}
\usepackage{amsmath, amssymb}
\usepackage{geometry}
\usepackage{setspace}
\usepackage{parskip}
\usepackage{enumitem}
\usepackage{hyperref}

\geometry{margin=2.5cm}
\onehalfspacing

\title{\textbf{Hito 2}}
\author{
    \textbf{Nombre Alumno:} Eduardo Caballero \\
    \textbf{Nombre Profesor guía:} Sebastián Cea
}
\date{}

\begin{document}

\maketitle

% ============================================================
\section*{Enfoque Metodológico}
% ============================================================

El presente Hito describe el enfoque metodológico propuesto para la memoria y especifica los
componentes preliminares del modelo a desarrollar. La estrategia consiste en partir del marco
de red dual desarrollado por Venegas (2024), quien a su vez se basa en el modelo de Jalan \&
Chakrabarti (2024) y lo extiende a un entorno con infraestructura tradicional y alternativa.
Sobre ese marco, la presente propuesta de memoria introduce heterogeneidad en aversión al
riesgo y en creencias sobre retornos como mecanismo para incorporar información asimétrica
entre los agentes. Esta extensión se conecta directamente con las oportunidades de mejora
identificadas en el Hito~1 y con las limitaciones que Venegas (2024) reconoce en la sección
5.3.3 de su tesis.

% ============================================================
\section{El marco de red dual}
% ============================================================

El modelo base sobre el cual se construye la memoria es el marco de red dual desarrollado por
Venegas (2024). Este marco extiende a Jalan \& Chakrabarti (2024) introduciendo dos
infraestructuras paralelas —una red Tradicional ($T$) y una red Alternativa ($A$)— sobre las
cuales tres agentes (Banco Nacional, Banco Local y Empresa Local) deciden simultáneamente
cómo asignar su exposición total $w_i^*$ entre ambas redes, sujetos a la restricción de
asignación:
\begin{equation}
    w_i^* = w_i^T + w_i^A.
\end{equation}

El marco se caracteriza formalmente por la tupla $(\mu_i,\, \gamma_i,\, \Sigma_i,\, \Psi_i^T,\,
\Psi_i^A,\, \tau,\, \lambda_A,\, w_i^*)$, donde $\mu_i$ es el vector de retornos esperados,
$\gamma_i$ el coeficiente de aversión al riesgo, $\Sigma_i$ la matriz de covarianzas, $\Psi_i^T$
y $\Psi_i^A$ las matrices de contrapartes permitidas en cada red, $\tau \geq 0$ los beneficios
de transparencia de la red alternativa, y $\lambda_A \geq 0$ los costos friccionales asociados a
su uso. La utilidad de cada agente toma la forma cuadrática estándar:
\begin{equation}
    g_i(W, P) = w_i^\top(\mu_i + \tau \cdot e_i^A - P \cdot e_i)
    - \gamma_i \cdot w_i^\top \Sigma_i\, w_i
    - \lambda_A \cdot \|w_i^A\|^2.
    \label{eq:utilidad}
\end{equation}

El resultado central de Venegas (2024) es que, bajo su calibración, la red dual recupera el
$100\%$ del bienestar perdido por la restricción regulatoria (mejora del $103{,}0\%$ sobre el
escenario restringido en el caso de tres agentes), mientras que la introducción de competencia
dentro de una única red cierra solo aproximadamente el $59\%$ de la brecha de bienestar
respecto al benchmark sin restricción. Este hallazgo establece la red alternativa como
mecanismo estructural de eficiencia y constituye el punto de partida sobre el cual se
construyen las preguntas de la presente memoria.

% ============================================================
\section{Limitación identificada y oportunidad de extensión}
% ============================================================

El resultado de Venegas (2024) se obtiene bajo dos supuestos simplificadores que el propio
autor reconoce como limitaciones en la sección 5.3.3 de su tesis. El primero, todos los agentes
comparten una aversión al riesgo idéntica ($\gamma_i = 1{,}0$ para todo $i$). El segundo, todos
los agentes mantienen creencias homogéneas sobre los retornos esperados ($\mu_i = \mu_j$).
Estos dos supuestos abstraen del marco precisamente las dos fuentes de heterogeneidad que,
en mercados reales, generan la mayor parte de las ganancias del comercio: la diferencia en
tolerancia al riesgo entre instituciones y la dispersión de creencias sobre la rentabilidad y
riesgo de la infraestructura alternativa.

Venegas (2024) explícitamente observa que la introducción de dispersión en creencias generaría
ganancias adicionales del comercio no capturadas por su marco, y que la heterogeneidad en
aversión al riesgo podría producir que agentes tolerantes al riesgo se concentren en la red
alternativa mientras que aquellos aversos al riesgo permanezcan en la red tradicional. El autor
concluye que sus estimaciones de recuperación del bienestar son, en consecuencia,
conservadoras. Esta observación abre directamente la oportunidad de extensión que motiva al
presente proyecto de memoria.

Esta oportunidad se alinea con las tres direcciones de mejora identificadas en el Hito~1. La
heterogeneidad en $(\gamma_i, \mu_i)$ es la palanca natural para incorporar información
asimétrica al modelo, porque permite que cada agente actúe sobre creencias y preferencias
propias, generando comportamiento estratégico diferenciado entre Banco Nacional, Banco
Local y Empresa Local.

% ============================================================
\section{Extensión propuesta y preguntas de investigación}
% ============================================================

La extensión propuesta consiste en relajar simultáneamente los dos supuestos de
homogeneidad identificados, permitiendo que cada agente (Banco Nacional, Banco Local,
Empresa Local) se caracterice por un coeficiente de aversión al riesgo $\gamma_i$ y un vector
de creencias sobre retornos $\mu_i$ propios, con $\gamma_i \neq \gamma_j$ y $\mu_i \neq \mu_j$
en general. El resto de la estructura del modelo de Venegas (2024) —como $\Psi_i^T$ y
$\Psi_i^A$, la transparencia $\tau$, los costos friccionales $\lambda_A$ y la restricción de
asignación total— se mantendría intacta. Se trata, en consecuencia, de una extensión
incremental que aísla el efecto de la heterogeneidad sobre los resultados de bienestar y
comportamiento estratégico, sin alterar el resto del andamiaje del marco base.

La asignación de perfiles a cada agente se propone con base en la lógica económica de sus
roles institucionales. El Banco Nacional, por su escala, se modelaría como un agente con
aversión al riesgo baja y creencias precisas sobre los retornos. El Banco Local opera con
aversión al riesgo intermedia y un grado parcial de información, reflejando su posición
intermedia entre el Banco Nacional y el sector real. Por su parte, la Empresa Local se
caracteriza por aversión al riesgo elevada y creencias ruidosas sobre los retornos de la
infraestructura alternativa, capturando la posición de un agente no-financiero con acceso
limitado a información especializada.

Esta extensión habilita tres preguntas de investigación que el marco homogéneo de Venegas
(2024) no puede formular:

\begin{enumerate}
    \item ?`Persiste la recuperación completa de bienestar cuando los agentes son heterogéneos,
    o la dispersión en $(\gamma_i, \mu_i)$ reduce la magnitud de la recuperación?

    \item ?`Se produce \textit{sorting} endógeno entre las dos redes? Es decir, ?`los agentes
    tolerantes al riesgo y bien informados (Banco Nacional) concentran su exposición en la red
    alternativa mientras los aversos al riesgo y peor informados (Empresa Local) permanecen
    en la red tradicional?

    \item ?`Cómo se distribuyen las ganancias de bienestar entre los tres agentes bajo
    heterogeneidad? ?`Obtiene el Banco Nacional un mayor porcentaje debido a su ventaja de
    acceso a información más profunda de los mercados?
\end{enumerate}

% ============================================================
\section*{Componentes Preliminares del Modelo}
% ============================================================

La especificación preserva la estructura del marco de Venegas (2024) y modifica únicamente
aquellos componentes directamente afectados por la heterogeneidad propuesta.

% ============================================================
\subsection{Agentes y parámetros heterogéneos}
% ============================================================

El modelo considera el conjunto de agentes $\mathcal{N} = \{1, 2, 3\}$, correspondientes al
Banco Nacional ($i = 1$), Banco Local ($i = 2$) y Empresa Local ($i = 3$). Cada agente se
caracteriza por dos parámetros individuales: un coeficiente de aversión al riesgo $\gamma_i > 0$
y un vector de creencias sobre retornos esperados $\mu_i \in \mathbb{R}^n$.

A diferencia de Venegas (2024), donde $\gamma_i = 1$ y $\mu_i = \mu$ para todo $i$, la
extensión propuesta requiere especificar una relación de orden cualitativa entre los parámetros
de los tres agentes. En términos de aversión al riesgo, se propone:
\begin{equation}
    \gamma_1 < \gamma_2 < \gamma_3,
\end{equation}
lo que refleja la jerarquía institucional descrita anteriormente. En términos de creencias, la
heterogeneidad se modela como dispersión alrededor de un valor de referencia $\bar{\mu}$
común, donde $\mu_1$ se aproxima a $\bar{\mu}$ con error pequeño, $\mu_2$ con error
intermedio, y $\mu_3$ con error grande. El resto de los parámetros individuales del marco base
—la matriz de covarianzas $\Sigma_i$ y la exposición total $w_i^*$— se conservan inicialmente
como en Venegas (2024).

% ============================================================
\subsection{Estructura de redes y restricciones}
% ============================================================

La estructura de redes se hereda íntegramente del marco de Venegas (2024). Cada agente $i$
opera simultáneamente sobre dos redes: la red Tradicional, caracterizada por la matriz
$\Psi_i^T$, y la red Alternativa, caracterizada por la matriz $\Psi_i^A$. Las matrices $\Psi$
señalan qué pares de agentes pueden mantener exposiciones entre sí en cada infraestructura.
Sobre esta estructura actúan dos parámetros globales que regulan qué tan ``atractiva'' es la
red alternativa: el premio de transparencia $\tau \geq 0$, que aumenta el retorno esperado de
las exposiciones tomadas en la red alternativa, y el costo friccional $\lambda_A \geq 0$, que
penaliza el tamaño de la posición tomada en dicha red.

Cada agente está sujeto a la restricción de asignación:
\begin{equation}
    w_i^* = w_i^T + w_i^A,
\end{equation}
que establece que la suma de las exposiciones tomadas en ambas redes debe igualar la
exposición total exógena del agente. Esta restricción captura la decisión central del modelo
—la asignación entre infraestructuras— y se mantiene sin modificación respecto al marco base.

\end{document}